% Options for packages loaded elsewhere
\PassOptionsToPackage{unicode}{hyperref}
\PassOptionsToPackage{hyphens}{url}
\documentclass[
  ignorenonframetext,
]{beamer}
\newif\ifbibliography
\usepackage{pgfpages}
\setbeamertemplate{caption}[numbered]
\setbeamertemplate{caption label separator}{: }
\setbeamercolor{caption name}{fg=normal text.fg}
\beamertemplatenavigationsymbolsempty
% remove section numbering
\setbeamertemplate{part page}{
  \centering
  \begin{beamercolorbox}[sep=16pt,center]{part title}
    \usebeamerfont{part title}\insertpart\par
  \end{beamercolorbox}
}
\setbeamertemplate{section page}{
  \centering
  \begin{beamercolorbox}[sep=12pt,center]{section title}
    \usebeamerfont{section title}\insertsection\par
  \end{beamercolorbox}
}
\setbeamertemplate{subsection page}{
  \centering
  \begin{beamercolorbox}[sep=8pt,center]{subsection title}
    \usebeamerfont{subsection title}\insertsubsection\par
  \end{beamercolorbox}
}
% Prevent slide breaks in the middle of a paragraph
\widowpenalties 1 10000
\raggedbottom
\AtBeginPart{
  \frame{\partpage}
}
\AtBeginSection{
  \ifbibliography
  \else
    \frame{\sectionpage}
  \fi
}
\AtBeginSubsection{
  \frame{\subsectionpage}
}
\usepackage{iftex}
\ifPDFTeX
  \usepackage[T1]{fontenc}
  \usepackage[utf8]{inputenc}
  \usepackage{textcomp} % provide euro and other symbols
\else % if luatex or xetex
  \usepackage{unicode-math} % this also loads fontspec
  \defaultfontfeatures{Scale=MatchLowercase}
  \defaultfontfeatures[\rmfamily]{Ligatures=TeX,Scale=1}
\fi
\usepackage{lmodern}
\ifPDFTeX\else
  % xetex/luatex font selection
\fi
% Use upquote if available, for straight quotes in verbatim environments
\IfFileExists{upquote.sty}{\usepackage{upquote}}{}
\IfFileExists{microtype.sty}{% use microtype if available
  \usepackage[]{microtype}
  \UseMicrotypeSet[protrusion]{basicmath} % disable protrusion for tt fonts
}{}
\makeatletter
\@ifundefined{KOMAClassName}{% if non-KOMA class
  \IfFileExists{parskip.sty}{%
    \usepackage{parskip}
  }{% else
    \setlength{\parindent}{0pt}
    \setlength{\parskip}{6pt plus 2pt minus 1pt}}
}{% if KOMA class
  \KOMAoptions{parskip=half}}
\makeatother
\usepackage{color}
\usepackage{fancyvrb}
\newcommand{\VerbBar}{|}
\newcommand{\VERB}{\Verb[commandchars=\\\{\}]}
\DefineVerbatimEnvironment{Highlighting}{Verbatim}{commandchars=\\\{\}}
% Add ',fontsize=\small' for more characters per line
\newenvironment{Shaded}{}{}
\newcommand{\AlertTok}[1]{\textcolor[rgb]{1.00,0.00,0.00}{\textbf{#1}}}
\newcommand{\AnnotationTok}[1]{\textcolor[rgb]{0.38,0.63,0.69}{\textbf{\textit{#1}}}}
\newcommand{\AttributeTok}[1]{\textcolor[rgb]{0.49,0.56,0.16}{#1}}
\newcommand{\BaseNTok}[1]{\textcolor[rgb]{0.25,0.63,0.44}{#1}}
\newcommand{\BuiltInTok}[1]{\textcolor[rgb]{0.00,0.50,0.00}{#1}}
\newcommand{\CharTok}[1]{\textcolor[rgb]{0.25,0.44,0.63}{#1}}
\newcommand{\CommentTok}[1]{\textcolor[rgb]{0.38,0.63,0.69}{\textit{#1}}}
\newcommand{\CommentVarTok}[1]{\textcolor[rgb]{0.38,0.63,0.69}{\textbf{\textit{#1}}}}
\newcommand{\ConstantTok}[1]{\textcolor[rgb]{0.53,0.00,0.00}{#1}}
\newcommand{\ControlFlowTok}[1]{\textcolor[rgb]{0.00,0.44,0.13}{\textbf{#1}}}
\newcommand{\DataTypeTok}[1]{\textcolor[rgb]{0.56,0.13,0.00}{#1}}
\newcommand{\DecValTok}[1]{\textcolor[rgb]{0.25,0.63,0.44}{#1}}
\newcommand{\DocumentationTok}[1]{\textcolor[rgb]{0.73,0.13,0.13}{\textit{#1}}}
\newcommand{\ErrorTok}[1]{\textcolor[rgb]{1.00,0.00,0.00}{\textbf{#1}}}
\newcommand{\ExtensionTok}[1]{#1}
\newcommand{\FloatTok}[1]{\textcolor[rgb]{0.25,0.63,0.44}{#1}}
\newcommand{\FunctionTok}[1]{\textcolor[rgb]{0.02,0.16,0.49}{#1}}
\newcommand{\ImportTok}[1]{\textcolor[rgb]{0.00,0.50,0.00}{\textbf{#1}}}
\newcommand{\InformationTok}[1]{\textcolor[rgb]{0.38,0.63,0.69}{\textbf{\textit{#1}}}}
\newcommand{\KeywordTok}[1]{\textcolor[rgb]{0.00,0.44,0.13}{\textbf{#1}}}
\newcommand{\NormalTok}[1]{#1}
\newcommand{\OperatorTok}[1]{\textcolor[rgb]{0.40,0.40,0.40}{#1}}
\newcommand{\OtherTok}[1]{\textcolor[rgb]{0.00,0.44,0.13}{#1}}
\newcommand{\PreprocessorTok}[1]{\textcolor[rgb]{0.74,0.48,0.00}{#1}}
\newcommand{\RegionMarkerTok}[1]{#1}
\newcommand{\SpecialCharTok}[1]{\textcolor[rgb]{0.25,0.44,0.63}{#1}}
\newcommand{\SpecialStringTok}[1]{\textcolor[rgb]{0.73,0.40,0.53}{#1}}
\newcommand{\StringTok}[1]{\textcolor[rgb]{0.25,0.44,0.63}{#1}}
\newcommand{\VariableTok}[1]{\textcolor[rgb]{0.10,0.09,0.49}{#1}}
\newcommand{\VerbatimStringTok}[1]{\textcolor[rgb]{0.25,0.44,0.63}{#1}}
\newcommand{\WarningTok}[1]{\textcolor[rgb]{0.38,0.63,0.69}{\textbf{\textit{#1}}}}
\setlength{\emergencystretch}{3em} % prevent overfull lines
\providecommand{\tightlist}{%
  \setlength{\itemsep}{0pt}\setlength{\parskip}{0pt}}
\usepackage{bookmark}
\IfFileExists{xurl.sty}{\usepackage{xurl}}{} % add URL line breaks if available
\urlstyle{same}
\hypersetup{
  hidelinks,
  pdfcreator={LaTeX via pandoc}}

\author{\texorpdfstring{}{}}
\date{}

\begin{document}

\section{What to record for reproducible
experiments}\label{what-to-record-for-reproducible-experiments}

\begin{frame}[fragile]{Minimum recording checklist}
\protect\phantomsection\label{minimum-recording-checklist}
The COGANT pipeline is deterministic on a fixed filesystem snapshot
under a fixed Python interpreter, so bit-for-bit reproduction of a
published run reduces to pinning exactly the inputs that the pipeline
consumes and the environment it consumes them in. For every experiment
whose output you intend to re-run, redistribute, or cite, record the
following:

\begin{enumerate}
\tightlist
\item
  \textbf{COGANT version and commit hash.} The \texttt{\_\_version\_\_}
  string in \texttt{../cogant/py/cogant/\_\_init\_\_.py} matches
  \texttt{project.version} in \texttt{../cogant/pyproject.toml}
  (currently \texttt{0.5.0}). Also capture the Git SHA of the checkout
  used --- \texttt{git\ rev-parse\ HEAD} from inside the
  \texttt{cogant/} subtree --- because a single version tag can span
  several bug-fix commits and the reproducibility contract is at commit
  granularity, not tag granularity. The shipped \texttt{cogant\ doctor}
  command prints both values.
\item
  \textbf{Python interpreter version.} Either \texttt{CPython\ 3.11.x},
  \texttt{3.12.x}, or \texttt{3.13.x} per the \texttt{classifiers} block
  in \texttt{pyproject.toml}; the canonical benchmark run used
  \texttt{CPython\ 3.12.11} on macOS arm64. Capture
  \texttt{python\ -\/-version} plus \texttt{uname\ -a} (or the
  equivalent on Windows / Linux) so that the OS and architecture are
  visible downstream.
\item
  \textbf{Dependency lock.} \texttt{uv\ sync\ -\/-extra\ all} resolves
  from \texttt{../cogant/uv.lock}; include the lockfile in any
  redistributed dataset so that optional extras (\texttt{viz},
  \texttt{multilang}, \texttt{all}) resolve to the exact versions that
  produced the published numbers. If the pinned
  \texttt{tree-sitter-javascript}, \texttt{tree-sitter-typescript},
  \texttt{pyarrow}, or \texttt{duckdb} versions drift, re-parsing the
  same JavaScript sources can produce different AST shapes and therefore
  different \texttt{program\_graph.json} outputs.
\item
  \textbf{Input repository commit hash.} For every target repository
  processed through the pipeline, record the Git commit hash of the
  snapshot at ingest time. The pipeline does not hash the input for you
  --- it simply walks the filesystem --- so an unpinned input is the
  most common source of non-reproducibility. For the shipped fixtures
  under \texttt{../cogant/examples/}, the relevant hash is the COGANT
  commit itself (the fixtures are vendored).
\item
  \textbf{Configuration file contents (redacted).} If you invoked the
  pipeline via \texttt{cogant.yaml} or \texttt{PipelineConfig(...)},
  persist the configuration next to the run outputs. Redact any
  credentials, private path components, or API tokens before
  publication; the config loader in \texttt{../cogant/py/cogant/config/}
  does not automatically strip these.
\item
  \textbf{List of stages executed.} The default ten-stage DAG is
  \texttt{ingest\ →\ static\ →\ normalize\ →\ graph\ →\ dynamic\ →\ translate\ →\ statespace\ →\ process\ →\ export\ →\ validate},
  but the CLI and \texttt{PipelineConfig} expose selective-execution
  flags (\texttt{-\/-no-dynamic}, \texttt{-\/-skip-validate},
  \texttt{PipelineConfig.skip\_stages}). Record which stages were
  enabled for the published run so that a reproducer does not silently
  re-enable an optional stage whose output was not part of the original
  report.
\item
  \textbf{Random seeds for learned consumers.} COGANT itself is
  deterministic --- it does not invoke any stochastic components in the
  default pipeline. However, any \textbf{downstream learned component}
  (for example an optional embedding model supplied to the GNN exporter,
  or a graph neural network trained on COGANT's exports) \emph{is}
  stochastic and must record its own seeds. Include the seed with the
  derived dataset, not with the COGANT bundle.
\item
  \textbf{Canonical output hashes.} After each run, hash the
  \texttt{gnn\_package/} directory and the top-level
  \texttt{bundle.json} with a content-addressable digest
  (\texttt{sha256} is sufficient). Store these hashes alongside the run
  outputs; republishing the same hash across different machines is the
  strongest reproducibility signal available to a reader.
\end{enumerate}
\end{frame}

\begin{frame}[fragile]{How to re-run the canonical fixture experiments}
\protect\phantomsection\label{how-to-re-run-the-canonical-fixture-experiments}
All numbers in Tables 4--7 of §6.3 and in the benchmark suite of §6.4
are regenerated by two scripts that live inside the COGANT package tree
and take no arguments beyond a \texttt{-\/-output} directory:

\begin{Shaded}
\begin{Highlighting}[]
\CommentTok{\# One{-}shot regeneration of every figure and the metrics.json summary}
\ExtensionTok{uv}\NormalTok{ run python ../cogant/evaluation/figures/generate\_figures.py }\DataTypeTok{\textbackslash{}}
    \AttributeTok{{-}{-}output}\NormalTok{ ../cogant/evaluation/figures/}
\end{Highlighting}
\end{Shaded}

\begin{Shaded}
\begin{Highlighting}[]
\CommentTok{\# Benchmark harness — three iterations per fixture, stage{-}level timing}
\ExtensionTok{uv}\NormalTok{ run python ../cogant/benchmarks/bench\_suite.py }\DataTypeTok{\textbackslash{}}
    \AttributeTok{{-}{-}iterations}\NormalTok{ 3 }\DataTypeTok{\textbackslash{}}
    \AttributeTok{{-}{-}output}\NormalTok{ ../cogant/benchmarks/results/}
\end{Highlighting}
\end{Shaded}

The figure-generation script re-runs the \texttt{RoundtripOrchestrator}
(\texttt{../cogant/examples/orchestrate\_roundtrip.py}) against every
fixture under \texttt{../cogant/examples/control\_positive/} and
\texttt{../cogant/examples/real\_world/}, re-reads the emitted JSON,
writes \texttt{metrics.json} alongside \texttt{fig1\_graph\_sizes.png},
\texttt{fig2\_node\_kinds.png}, \texttt{fig3\_state\_space.png}, and
\texttt{fig4\_pipeline\_latency.png}, and finally copies the canonical
figure set back into the \texttt{../cogant/evaluation/figures/} tree.
The benchmark-suite harness times the bare translation pipeline
(\texttt{ingest}, \texttt{static}, \texttt{normalize}, \texttt{graph},
\texttt{translate}, \texttt{statespace}) and writes a dated Markdown
report of the form \texttt{suite\_YYYYMMDD.md} alongside a
machine-readable JSON sidecar; the report committed as
\texttt{../cogant/benchmarks/results/suite\_20260409.md} is the
canonical source for Table 11 in §6.4.
\end{frame}

\begin{frame}[fragile]{What the manifest.json index records}
\protect\phantomsection\label{what-the-manifest.json-index-records}
Every \texttt{gnn\_package/} directory emitted by the pipeline includes
a \texttt{manifest.json} file whose job is to close the reproducibility
loop without requiring the downstream consumer to re-hash the bundle.
The manifest records: the COGANT version that produced the bundle, the
interpreter and platform identifiers from step 2 of the checklist, the
list of stages actually executed (step 6), the schema name passed to the
state-space compiler, the paths of every file in the
\texttt{gnn\_package/} directory, and the SHA-256 hash of each file. A
reproducer can therefore verify a published bundle with a single pass
over the manifest, without consulting any external metadata file.
\end{frame}

\begin{frame}[fragile]{Worked example: regenerating the Table 4 row for
\texttt{flask\_app}}
\protect\phantomsection\label{worked-example-regenerating-the-table-4-row-for-flask_app}
The canonical Table 4 row for the \texttt{flask\_app} fixture (six
files, 853 lines, 98 nodes, 597 edges, 14.58 s wall-clock) was produced
as follows, and the command sequence below should bit-for-bit reproduce
the same \texttt{program\_graph.json} and \texttt{state\_space.json} on
any macOS arm64 machine with the same uv lockfile:

\begin{Shaded}
\begin{Highlighting}[]
\BuiltInTok{cd}\NormalTok{ projects\_in\_progress/cogant/cogant           }\CommentTok{\# package root}
\ExtensionTok{uv}\NormalTok{ sync }\AttributeTok{{-}{-}extra}\NormalTok{ all                               }\CommentTok{\# pins every dep per uv.lock}
\FunctionTok{git}\NormalTok{ rev{-}parse HEAD                                }\CommentTok{\# record the commit hash}
\ExtensionTok{python} \AttributeTok{{-}{-}version} \KeywordTok{\&\&} \FunctionTok{uname} \AttributeTok{{-}a}                      \CommentTok{\# record interpreter + platform}
\ExtensionTok{uv}\NormalTok{ run cogant translate }\DataTypeTok{\textbackslash{}}
\NormalTok{    examples/real\_world/flask\_app }\DataTypeTok{\textbackslash{}}
    \AttributeTok{{-}{-}output}\NormalTok{ /tmp/repro\_flask\_app }\DataTypeTok{\textbackslash{}}
    \AttributeTok{{-}{-}layout{-}output}
\ExtensionTok{uv}\NormalTok{ run cogant validate /tmp/repro\_flask\_app/gnn\_package}
\FunctionTok{sha256sum}\NormalTok{ /tmp/repro\_flask\_app/gnn\_package/}\PreprocessorTok{*}\NormalTok{.json }\KeywordTok{|} \FunctionTok{sort}
\end{Highlighting}
\end{Shaded}

The \texttt{validate} step should report score \texttt{100.0\ /\ 100}
with zero errors and zero warnings, and the sorted SHA-256 listing
should match the one recorded in
\texttt{../cogant/examples/real\_world/flask\_app/output/manifest.json}
(committed alongside the fixture). Any mismatch points to either a
dependency drift (step 3) or an input drift (the Flask fixture itself
changed between the committed output and the re-run).
\end{frame}

\begin{frame}[fragile]{Data ethics and licensing for exported bundles}
\protect\phantomsection\label{data-ethics-and-licensing-for-exported-bundles}
Every COGANT bundle can contain identifiers, docstrings, and inline
comments lifted verbatim from the input source code. When redistributing
a bundle derived from a third-party repository --- for example when
publishing a dataset of \texttt{gnn\_package/} outputs from open-source
Python projects --- the downstream license terms must be respected: the
original repository's license applies to any prose or identifier that
the bundle quotes, and the COGANT MIT license covers only the pipeline
code and the synthetic fixtures it ships. Organisations with private
data policies should additionally review bundles for personally
identifiable information before publication; the pipeline does not
perform PII scrubbing on its own. The short policy statement is in §7
(see \textbf{Data ethics and licensing} in
\href{07_reproducibility.md}{\texttt{07\_reproducibility.md}}); this
subsection supplies the actionable hashing and manifest-recording
recipe.
\end{frame}

\begin{frame}[fragile]{Cross-references}
\protect\phantomsection\label{cross-references}
\begin{itemize}
\tightlist
\item
  The CLI hub at
  \href{../cogant/docs/cli/README.md}{\texttt{../cogant/docs/cli/README.md}}
  links to every flag that changes the recorded-output shape, and the
  stage list in
  \href{../cogant/docs/architecture/README.md}{\texttt{../cogant/docs/architecture/README.md}}
  enumerates the ten-stage DAG.
\item
  The per-release narrative in \texttt{../cogant/CHANGELOG.md} and
  \texttt{../cogant/RELEASE\_NOTES.md} documents which default-behaviour
  changes (for example the v0.5.0 POLICY / CONTEXT stub-emission fix
  discussed in §5 and \texttt{ROUNDTRIP\_IMPROVEMENT.md}) could affect a
  re-run against a pre-v0.5.0 bundle.
\item
  The calibration sweep plan in
  \texttt{../cogant/docs/evaluation/CALIBRATION.md} is the canonical
  reference for the \texttt{TODO(calibration)} markers cited in §5;
  re-running a confidence-threshold sweep requires the 20+ repository
  gold-standard corpus discussed there.
\end{itemize}
\end{frame}

\end{document}
